\documentclass[aspectratio=169]{beamer}

\usepackage[T2A]{fontenc}
\usepackage[utf8]{inputenc}
\usepackage[bulgarian]{babel}
\usepackage{tempora}
\usepackage{amsmath}
\usepackage{graphicx}
\usepackage{booktabs}
\usepackage{array}
\usepackage{tikz}
\usepackage{pgfplots}
\pgfplotsset{compat=1.18}

\usetheme{Copenhagen}
\usecolortheme{default}
\setbeamertemplate{navigation symbols}{}

\title{Моделиране на център за обработка на заявки с два компютъра}
\author{Кирил Вълков \\ фак. № 121222194}
\institute{Технически университет -- София \\ Факултет „Компютърни системи и технологии“}
\date{}

\begin{document}

\begin{frame}
    \titlepage
\end{frame}

\begin{frame}{Съдържание}
    \tableofcontents
\end{frame}

\section{Постановка на задачата}

\begin{frame}{Постановка на задачата}
    \begin{itemize}
        \item Центърът за обработка на заявки се състои от два компютъра.
        \item Към Компютър 1 заявките пристигат с експоненциално разпределение и средно междупристигателно време 1 s.
        \item Към Компютър 2 заявките пристигат с нормално разпределение и средно междупристигателно време 3 s с отклонение в рамките на $\pm 1$ s.
        \item Времето за обслужване и за двата компютъра е $(2 \pm 1)$ s, тоест равномерно в интервала $[1,3]$ s.
        \item Периодът на моделиране е 1800 s (30 min).
    \end{itemize}
\end{frame}

\section{Теоретичен анализ}

\begin{frame}{Теоретичен анализ на системата}
    Системата може да се разглежда като две независими системи за масово обслужване.

    \vspace{0.5cm}

    \begin{block}{Компютър 1}
        M/G/1 система с $E[A]=1$ s и $E[S]=2$ s.
    \end{block}

    \begin{block}{Компютър 2}
        G/G/1 система с приблизително $E[A]=3$ s и $E[S]=2$ s.
    \end{block}
\end{frame}

\begin{frame}{Коефициент на натоварване}
    \begin{columns}[T]
        \begin{column}{0.48\textwidth}
            \begin{block}{Компютър 1}
                \[
                    \rho_1 = \frac{E[S]}{E[A]} = \frac{2}{1} = 2
                \]
                \begin{itemize}
                    \item $\rho_1 > 1$
                    \item Системата е нестабилна
                    \item Опашката нараства неограничено
                \end{itemize}
            \end{block}
        \end{column}
        \begin{column}{0.48\textwidth}
            \begin{block}{Компютър 2}
                \[
                    \rho_2 = \frac{2}{3} \approx 0.67
                \]
                \begin{itemize}
                    \item $\rho_2 < 1$
                    \item Системата е стабилна
                    \item Изчакването е минимално
                \end{itemize}
            \end{block}
        \end{column}
    \end{columns}
\end{frame}

\section{Реализация на модела}

\begin{frame}{Реализация на симулацията}
    \begin{itemize}
        \item Основната симулация е реализирана на Python с библиотеката \texttt{simpy}.
        \item Използвани са два независими генератора на заявки.
        \item Всеки компютър е моделиран като ресурс с капацитет 1.
        \item За валидация е създаден и алтернативен модел на GPSS.
    \end{itemize}

    \vspace{0.4cm}

    \begin{block}{Основни параметри}
        \begin{itemize}
            \item Време за симулация: 1800 s
            \item Компютър 1: експоненциални пристигания, средно 1 s
            \item Компютър 2: нормални пристигания, средно 3 s
            \item Обслужване: равномерно в $[1,3]$ s
        \end{itemize}
    \end{block}
\end{frame}

\begin{frame}{Логика на модела}
    \begin{enumerate}
        \item Генерират се заявки за двата компютъра според зададените разпределения.
        \item Всяка заявка постъпва в опашка към съответния ресурс.
        \item При освобождаване на ресурса заявката започва обслужване.
        \item Измерват се:
        \begin{itemize}
            \item броят обработени заявки,
            \item средно време на изчакване,
            \item максимално време на изчакване,
            \item натовареност на ресурса.
        \end{itemize}
        \item Получените резултати се сравняват между Python и GPSS.
    \end{enumerate}
\end{frame}

\begin{frame}[fragile]{GPSS модел на системата}
    \small
    \begin{block}{Основен GPSS код}
\scriptsize
\begin{verbatim}
SIMULATE
GENERATE 1
SEIZE 1
ADVANCE 2,1
RELEASE 1
TERMINATE 0
GENERATE 3,1
SEIZE 2
ADVANCE 2,1
RELEASE 2
TERMINATE 0
GENERATE 1800
TERMINATE 1
START 1
\end{verbatim}
    \end{block}
\end{frame}

\begin{frame}{Резултати от GPSS симулацията}
    \small
    \begin{table}
        \centering
        \begin{tabular}{lcc}
            \toprule
            \textbf{Параметър} & \textbf{Компютър 1} & \textbf{Компютър 2} \\
            \midrule
            Входове във FACILITY & 910 & 595 \\
            Средно време за обслужване (s) & 1.977 & 2.002 \\
            Натовареност & 0.999 & 0.662 \\
            Забавени заявки & 890 & 0 \\
            \bottomrule
        \end{tabular}
    \end{table}

    \vspace{0.3cm}

    \begin{itemize}
        \item GPSS отчетът показва силно претоварване на Компютър 1 и липса на натрупване при Компютър 2.
        \item Получените стойности са в съответствие с Python симулацията и теоретичния анализ.
    \end{itemize}
\end{frame}

\section{Резултати}

\begin{frame}{Резултати от симулацията}
    \begin{table}
        \centering
        \begin{tabular}{>{\raggedright\arraybackslash}p{4.4cm}cc}
            \toprule
            \textbf{Параметър} & \textbf{Компютър 1} & \textbf{Компютър 2} \\
            \midrule
            Обработени заявки & 895 & 597 \\
            Средно време на изчакване (s) & 444.78 & 0.016 \\
            Максимално време на изчакване (s) & 932.10 & 0.91 \\
            Натовареност & 99.93\% & 67.27\% \\
            \bottomrule
        \end{tabular}
    \end{table}
\end{frame}

\begin{frame}{Сравнение между Python и GPSS}
    \begin{table}
        \centering
        \begin{tabular}{lcccc}
            \toprule
            \textbf{Параметър} & \multicolumn{2}{c}{\textbf{Компютър 1}} & \multicolumn{2}{c}{\textbf{Компютър 2}} \\
            \cmidrule(lr){2-3} \cmidrule(lr){4-5}
             & \textbf{Python} & \textbf{GPSS} & \textbf{Python} & \textbf{GPSS} \\
            \midrule
            Обработени заявки & 895 & 910 & 597 & 595 \\
            Натовареност & 99.93\% & 99.9\% & 67.27\% & 66.2\% \\
            \bottomrule
        \end{tabular}
    \end{table}

    \vspace{0.3cm}

    \begin{itemize}
        \item Резултатите са близки и потвърждават коректността на модела.
        \item Разликите се дължат на случайния характер и различните реализации на разпределенията.
    \end{itemize}
\end{frame}

\begin{frame}{Графична интерпретация на изчакването}
    \centering
    \begin{tikzpicture}
        \begin{axis}[
            width=11.5cm,
            height=6.3cm,
            xlabel={Време на изчакване (s)},
            ylabel={Брой заявки},
            title={Хистограма на времената за изчакване за Компютър 1},
            ybar,
            ymin=0,
            xmin=0,
            xmax=1000,
            bar width=6pt
        ]
        \addplot coordinates {
            (46.6, 86) (139.8, 118) (233.0, 100) (326.2, 93) (419.4, 92)
            (512.7, 75) (605.9, 88) (699.1, 80) (792.3, 80) (885.5, 84)
        };
        \end{axis}
    \end{tikzpicture}
\end{frame}

\begin{frame}{Основни изводи от резултатите}
    \begin{itemize}
        \item Компютър 1 работи в претоварен режим.
        \item Натовареността му е практически 100\%, а опашката расте с времето.
        \item Компютър 2 остава в стабилен режим с почти нулево време на изчакване.
        \item Теоретичният анализ и симулационните резултати съвпадат.
    \end{itemize}
\end{frame}

\section{Заключение}

\begin{frame}{Заключение}
    \begin{block}{Обобщение}
        Симулационният модел показва, че при зададените параметри първият компютър не може да обслужва входящия поток без натрупване на опашка, докато вторият компютър работи устойчиво.
    \end{block}

    \vspace{0.4cm}

    \begin{alertblock}{Препоръка}
        За оптимизация на системата е целесъобразно част от натоварването на Компютър 1 да бъде пренасочено към Компютър 2 или да се увеличи производителността на първия компютър.
    \end{alertblock}
\end{frame}

\begin{frame}
    \centering
    \Huge Благодаря за вниманието!
\end{frame}

\end{document}
